\chapter{Fehlersuche}
\label{kap:fehlersuche}

Dieses Kapitel sammelt die häufigsten Meldungen und Probleme mit ihren
Ursachen. Meldungen des Simulators sind englisch; im Editor stehen sie im
\menue{Protokoll} bzw. in \menue{Probleme}.

\section{Fehler in der Datei}

\begin{description}[style=nextline,leftmargin=2em]
  \item[{\code{error: JSON syntax error: ... parse error at line 12, column 5}}]
    Ein Komma fehlt oder ist zu viel (nach dem letzten Element einer Liste
    ist keines erlaubt), eine Klammer oder ein Anführungszeichen ist offen.
    Der Editor markiert die Zeile; oft liegt der Fehler in der Zeile
    davor.
  \item[{\code{warning: components[0]: unknown parameter "valu" ignored}}]
    Ein Tippfehler in einem Schlüssel. Der Wert wird übergangen, die
    Rechnung läuft weiter – mit dem Standardwert. Warnungen deshalb ernst
    nehmen.
  \item[{\code{error: component 'D1': model "XYZ" is not defined}}]
    Das Modell fehlt unter \code{models}. Im Editor das Modell über die
    Bibliothek übernehmen oder den Namen korrigieren.
  \item[{\code{error: components[0]: expected nodes [a, b]}}]
    Falsche Zahl von Knoten für diesen Bauteiltyp (Tabelle in
    Kapitel~\ref{kap:format-bauteile}).
  \item[{\code{error: output "v(x)": unknown node 'x'}}]
    Ein Ausgang nennt einen Knoten, den es nicht gibt – oft nach dem
    Umbenennen eines Labels im Plan, oder ein Knoten in einer
    Teilschaltung ohne \code{Instanz.}-Präfix.
  \item[{\code{error: circuit has no ground node}}]
    Kein Knoten heißt \code{"0"} oder \code{"gnd"}. Im Plan ein
    Massesymbol setzen (\taste{G}).
\end{description}

\section{Der Arbeitspunkt konvergiert nicht}
\label{sec:schwebend}

\begin{shell}
error: DC operating point did not converge (check for floating nodes,
       loops of voltage sources/inductors, or try a nodeset)
\end{shell}

Die häufigsten Ursachen:

\begin{description}[style=nextline,leftmargin=2em]
  \item[Schwebende Knoten]
    Ein Knoten ohne Gleichstrompfad nach Masse – etwa zwischen zwei
    Kondensatoren, an einer Stromquelle mit Kondensator oder auf der
    Sekundärseite eines Übertragers. Seine Spannung ist im Arbeitspunkt
    unbestimmt. Abhilfe: ein großer Widerstand (z.\,B. 100\,M\ohm) nach
    Masse. Manchmal findet esim trotzdem eine Lösung; die Spannung des
    Knotens ist dann willkürlich (oft 0\,V), Differenzen wie
    \code{v(plus,minus)} sind davon nicht betroffen.
  \item[Maschen aus Spannungsquellen oder Spulen]
    Zwei parallele Spannungsquellen verschiedener Spannung, eine Spule
    direkt parallel zu einer Spannungsquelle: Im Arbeitspunkt (Spule =
    Kurzschluss) widersprechen sich die Gleichungen. Abhilfe: ein kleiner
    Serienwiderstand.
  \item[Stromquelle in Reihe mit einem Kondensator oder offenen Zweig]
    Der Strom kann nicht fließen. Abhilfe: Parallelwiderstand.
  \item[Schaltungen mit Mitkopplung]
    Flipflops, Schmitt-Trigger, Oszillatoren haben mehrere oder keine
    stabilen Arbeitspunkte. \code{"nodeset"} gibt einen Startwert vor; bei
    Oszillatoren hilft oft \code{"uic": true} in der Transientenanalyse.
  \item[Extreme Werte]
    Widerstände von Femtoohm oder Teraohm, Verstärkungen über $10^{9}$
    verschlechtern die Kondition. Werte prüfen (\code{M} statt \code{m}?).
\end{description}

Bevor esim aufgibt, versucht es nacheinander Newton, eine
Pseudo-Transiente, gmin-Stepping und Quellen-Stepping. Wenn keines zum Ziel
führt, liegt fast immer einer der obigen Gründe vor.

\section{Probleme in der Transientenanalyse}

\begin{description}[style=nextline,leftmargin=2em]
  \item[Rechnung sehr langsam]
    Die Statistik (\code{steps}, \code{rejected}) zeigt, ob viele Schritte
    verworfen werden. Ursachen sind meist sehr schnelle Vorgänge (steile
    Flanken, kleine Kapazitäten an großen Widerständen), Schalter ohne
    Hysterese oder ein zu kleines \code{tstep}. \code{lte_floor: 1} hilft bei
    schwingenden Signalen; für Schaltnetzteile, die lange einschwingen, ist
    \code{pss} oft viel schneller.
  \item[{\code{more than 100000 events ... (chattering)}}]
    Ein Vergleich, Schalter oder eine ideale Diode schaltet ohne Ende hin
    und her, z.\,B. \code{der(x) = if(x < 1, 1, -1)}. Abhilfe: Hysterese
    (\code{vh} am Schalter, zwei Ereignisse mit verschiedenen Schwellen)
    oder eine geglättete Kennlinie (\code{tanh} statt \code{if}).
  \item[Kurven eckig oder Spitzen fehlen]
    Die Kurve hat zu wenige Punkte: \code{tstep} kleiner oder
    \code{output_step} setzen.
  \item[Schwingung klingt künstlich ab]
    Numerische Dämpfung von BDF: \code{"method": "trap"}.
  \item[Strom einer Spannungsquelle zappelt]
    Typisch für \code{"method": "trap"} nach Sprüngen; \code{bdf} oder
    \code{gear2} verwenden.
  \item[Anfangszustand nicht wie erwartet]
    Ohne \code{"uic": true} startet die Rechnung im Arbeitspunkt; \code{ic}
    an Kondensatoren gilt nur mit \code{uic}.
\end{description}

\section{Probleme im Editor}

\begin{description}[style=nextline,leftmargin=2em]
  \item[Der Plan lässt sich nicht ändern]
    Der Text ist kein gültiges JSON. Der Plan zeigt den letzten gültigen
    Stand, bis der Fehler (siehe \menue{Probleme}) behoben ist.
  \item[Roter Ring an einem Anschluss]
    Plan und Text verbinden den Anschluss verschieden
    (Kapitel~\ref{kap:editor-schaltplan}). Gerechnet wird der Text; die
    nächste Änderung im Plan übernimmt die Knoten des Plans.
  \item[\tasten{Alt}{Klick} zeigt keine Kurve]
    Die Größe ist nicht gespeichert: unter \code{"outputs"} eintragen (die
    Statuszeile nennt den Namen) und neu simulieren.
  \item[»Der Simulator ist abgestürzt« oder esim nicht gefunden]
    \esim\ muss neben dem Editor liegen (im Paket mitgeliefert). Beim
    selbst gebauten Editor kopiert der Build es dorthin.
  \item[Englische Dialoge unter Windows]
    Im Programmordner fehlen die Übersetzungen
    (\datei{translations/qt_de.qm}); das Paket aus \code{cmake --install}
    bringt sie mit.
\end{description}

\section{Wenn nichts hilft}

\begin{itemize}
  \item Die Schaltung verkleinern, bis der Fehler verschwindet: Teile
    auskommentieren (\tasten{Strg}{/}), Analysen einzeln rechnen.
  \item Mit \code{op} beginnen; erst wenn der Arbeitspunkt plausibel ist,
    die anderen Analysen.
  \item Das Protokoll ganz lesen: Warnungen weiter oben erklären oft den
    Fehler am Ende.
\end{itemize}
